\documentclass[a4paper,12pt]{article}
\usepackage{docsTemplate}
\usepackage{longtable}
\usepackage{float}
\usepackage[official]{eurosym}
\usepackage{pgf-pie}
\usepackage{makecell}

\usepackage{titlesec}

\titleformat{\paragraph}[block]
  {\normalfont\normalsize\bfseries}
  {\theparagraph}{1em}{}

\titlespacing*{\paragraph}{0pt}{3.25ex plus 1ex minus .2ex}{1.5ex}

\setTitle{Piano di Progetto}
\setAuthors{Bianca Zaghetto}
\setVerificators{Mattia Oliva
Medin \\ & Sara Gioia
Fichera}
\setApprovation{}
\setVersion{0.3}
\setType{Esterno}
\setDestination{Tullio Vardanega \\ & Riccardo Cardin \\ & Sanmarco Informatica \\[-0.1em] & Team Three Technologies}
\setcounter{secnumdepth}{4}
\setcounter{tocdepth}{4}

\begin{document}
\include{docTitlepage}

\newpage
\section*{Registro delle modifiche}

\begin{table}[h!]
\rowcolors{2}{lightgray}{white}
\begin{tabularx}{\textwidth}{|l|l|l|l|L|L|}
\hline
\textbf{Vers.} & \textbf{Data} & \textbf{Autore} & \textbf{Ruolo} & \textbf{Verifica} & \textbf{Oggetto} \\
\hline
\textbf{0.3} & 09/12/25 & Bianca Zaghetto & Verificatore & Sara Gioia Fichera & Sezione calendario \\
\hline
\textbf{0.2} & 02/12/25 & Bianca Zaghetto & Verificatore & Sara Gioia Fichera & Sezione introduttiva e sezione rischi attesi \\
\hline
\textbf{0.1} & 27/11/25 & Bianca Zaghetto & Verificatore & Mattia Oliva Medin & Prima stesura: struttura documento \\
\hline
\end{tabularx}
\end{table}

\newpage
\tableofcontents

\newpage
\listoffigures

\newpage
\listoftables

\newpage

\section{Introduzione}
    \subsection{Scopo del documento}
        Il presente documento ha lo scopo di tenere traccia e definire le modalità di svolgimento delle attività nonché la suddivisione delle risorse in esse impiegate.
        Per la stesura del Piano di Progetto viene applicato un approccio incrementale, integrando informazioni fino al termine del progetto. Il documento tratta i seguenti temi:
        \begin{itemize}
            \item Rischi attesi e loro mitigazione
            \item Calendario di massima del progetto
            \item Impegno orario previsto e stima dei costi totali
            \item Periodi di avanzamento
            \item Consuntivo di periodo
            \item Preventivo
        \end{itemize}
    \subsection{Glossario}
        Per definizioni esaustive dei termini utilizzati in questo documento, si rimanda al Glossario disponibile nella documentazione. Tutti i termini che si possono trovare in esso vengono identificati da una G a pedice.
    \subsection{Riferimenti}
        \subsubsection{Riferimenti normativi}
            \begin{itemize}
                \item \href{https://team-three-technologies.github.io/Documenti/1%20-%20RTB/Norme_di_progetto_v0_4.pdf}{Norme di Progetto} (Ultimo accesso: 11/01/2026)
                \item \href{https://www.math.unipd.it/~tullio/IS-1/2025/Dispense/PD1.pdf}{Regolamento del progetto didattico} (Ultimo accesso: 11/01/2026)
            \end{itemize}
        \subsubsection{Riferimenti informativi}
            \begin{itemize}
                \item \href{https://team-three-technologies.github.io/Documenti/1%20-%20RTB/Analisi_dei_Requisiti_v0_1.pdf}{Analisi dei Requisiti} (Ultimo accesso: 08/01/2026)
                \item Verbali interni
                \item Verbali esterni
                \item \href{https://www.math.unipd.it/~tullio/IS-1/2025/Progetto/C3.pdf}{Capitolato C3} (Ultimo accesso: 10/01/2026)
                \item \href{https://www.math.unipd.it/~tullio/IS-1/2025/Dispense/T02.pdf}{Processi di ciclo di vita del \textit{software}} (Ultimo accesso: 05/01/2026)
                \item \href{https://www.math.unipd.it/~tullio/IS-1/2025/Dispense/T04.pdf}{Gestione di progetto} (Ultimo accesso: 07/01/2026)
            \end{itemize}

\newpage
\section{Rischi attesi e loro mitigazione}
\label{sec:rischi}
    \subsection{Introduzione}
        In questa sezione vengono elencati i rischi che si potrebbero incontrare durante lo svolgimento del progetto, corredati dalla possibile mitigazione relativa.
        
        Per rendere facilmente fruibile il documento, viene associato un codice identificativo univoco ad ogni rischio.
    \subsection{R1 - Tecnologie di lavoro e di produzione SW}
        \begin{table}[H]
        \rowcolors{2}{lightgray}{white}
        \label{tab:r1} 
            \begin{tabularx}{\textwidth}{|l|L|L|L|}
            \hline
            \textbf{Codice} & \textbf{Oggetto} & \textbf{Descrizione} & \textbf{Mitigazione}\\
            \hline
            \textbf{R1-1} & Inesperienza nelle tecnologie da utilizzare & Uno o più membri del gruppo si trovano in difficoltà nell'utilizzo delle tecnologie scelte & Studio personale e condivisione delle conoscenze\\
            \hline
            \textbf{R1-2} & Problemi tecnici software o hardware & Uno o più membri del gruppo riscontrano problemi tecnici con strumenti software o hardware & Si provvederà ad aggiornare il gruppo sul problema riscontrato. Qualora il membro o i membri siano impossibilitati a portare a termine un obiettivo, dovranno trovare una soluzione alternativa o cedere l'incarico a qualcun altro finché non avranno risolto il problema.\\            
            \hline
            \end{tabularx}
            \caption{Rischi sulle tecnologie di lavoro e di produzione SW}
        \end{table}
        
    \subsection{R2 - Rapporti interpersonali}
        \begin{table}[H]
        \rowcolors{2}{lightgray}{white}
        \label{tab:r2}
            \begin{tabularx}{\textwidth}{|l|L|L|L|}
            \hline
            \textbf{Codice} & \textbf{Oggetto} & \textbf{Descrizione} & \textbf{Mitigazione}\\
            \hline
            \textbf{R2-1} & Divergenze di pensiero & Il gruppo non riesce a giungere ad una decisione unanime per divergenze di pensiero & Utilizzo di sondaggi gestiti dal responsabile o, in casi estremi, richiesta di intervento da parte del docente/proponente\\
            \hline
            \textbf{R2-2} & Mancanza di riscontro & Uno o più membri del gruppo non danno riscontro rispetto agli incarichi assegnati & Promemoria regolari, se necessario per canali privati o in presenza. In casi estremi verrà contattato il docente.\\
            \hline
            \textbf{R2-3} & Abbandono del progetto & Uno o più membri del gruppo decidono di abbandonare il progetto per motivi personali & Verrà contattato il docente per rendere nota la situazione e si provvederà, se necessario, a valutare insieme alla proponente una negoziazione dei requisiti richiesti \\
            \hline
            \end{tabularx}
            \caption{Rischi sui rapporti interpersonali}
        \end{table}
        
    \subsection{R3 - Organizzazione del lavoro}
        \begin{table}[H]
        \rowcolors{2}{lightgray}{white}
        \label{tab:r3}
            \begin{tabularx}{\textwidth}{|l|L|L|L|}
            \hline
            \textbf{Codice} & \textbf{Oggetto} & \textbf{Descrizione} & \textbf{Mitigazione}\\
            \hline
            \textbf{R3-1} & Assenze alle riunioni esterne/interne & Uno o più membri del gruppo non presenziano alle riunioni interne/esterne & Condivisione appunti e registrazioni delle riunioni per permettere il recupero di quanto svolto\\
            \hline
            \textbf{R3-2} & Isolamento & Un membro si isola dal gruppo, non risponde alle comunicazioni, non presenzia costantemente alle riunioni e non svolge gli incarichi previsti & Confronto diretto sulle cause dell'isolamento, analisi delle possibili soluzioni ed eventuale riorganizzazione. In casi gravi, si provvederà a contattare il docente per decidere come procedere\\
            \hline
            \end{tabularx}
            \caption{Rischi sull'organizzazione del lavoro}
        \end{table}
        
    \subsection{R4 - Requisiti e rapporti con gli stakeholder}
        \begin{table}[H]
        \rowcolors{2}{lightgray}{white}
        \label{tab:r4}
            \begin{tabularx}{\textwidth}{|l|L|L|L|}
            \hline
            \textbf{Codice} & \textbf{Oggetto} & \textbf{Descrizione} & \textbf{Mitigazione}\\
            \hline
            \textbf{R4-1} & Malcomprensione requisiti & Il gruppo non comprende appieno i requisiti richiesti o li comprende in modo sbagliato & Utilizzo del form per la comunicazione asincrona con la proponente, ponendo dubbi e domande. Analisi dettagliata dei requisiti con la proponente nel corso delle riunioni. In casi eccezionali, richiesta di riunioni aggiuntive.\\
            \hline
            \end{tabularx}
            \caption{Rischi sui requisiti e rapporti con gli stakeholder}
        \end{table}
        
    \subsection{R5 - Tempi e costi}
        \begin{table}[H]
        \rowcolors{2}{lightgray}{white}
        \label{tab:r5}
            \begin{tabularx}{\textwidth}{|l|L|L|L|}
            \hline
            \textbf{Codice} & \textbf{Oggetto} & \textbf{Descrizione} & \textbf{Mitigazione}\\
            \hline
            \textbf{R5-1} & Ritardo di consegna & Uno o più membri del gruppo non riescono a completare una consegna nei tempi previsti & Qualora possibile, si concederà del tempo aggiuntivo per completare la consegna. Altrimenti si procederà ad affiancare un membro del gruppo disponibile o, in casi estremi, ad effettuare una sostituzione\\
            \hline
            \textbf{R5-2} & Superamento budget & Per il completamento del progetto è richiesto un budget superiore a quello preventivato & Negoziazione al ribasso sui requisiti del progetto con il proponente \\
            \hline
            \end{tabularx}
            \caption{Rischi su tempi e costi}
        \end{table}

\newpage
\section{Calendario di massima del progetto}
    \subsection{Introduzione}
        In questa sezione vengono esposte le date previste per le revisioni del progetto, considerando quanto visto in \href{sec:rischi}{\textit{Rischi attesi e loro mitigazione}} e \href{sec:pianificazione}{\textit{Pianificazione}}.
    \subsection{Stesura iniziale}
        Gli obiettivi iniziali prefissati dal gruppo, per quanto concerne il calendario, sono i seguenti:
        \begin{table}[H]
            \centering
            \label{tab:cal1}
            \begin{tabular}{|c|c|}
            \hline
                \rowcolor{lightgray}{\textbf{Revisione}} & {\textbf{Data}}  \\
            \hline
                Requirements and Technology Baseline & 19/01/2026 \\
            \hline
                Product Baseline & 27/03/2026 \\
            \hline
            \end{tabular}
            \caption{Previsione iniziale delle date di consegna delle revisioni}
        \end{table}

    % \subsection{Prima Revisione}
    %     \begin{table}[H]
    %         \centering
    %         \begin{tabular}{|c|c|}
    %         \hline
    %             \rowcolor{lightgray}{\textbf{Revisione}} & {\textbf{Data}}  \\
    %         \hline
    %             Product Baseline & ???\\
    %         \hline
    %         \end{tabular}
    %         \caption{Data di consegna prevista aggiornata a seguito della RTB}
    %         \label{tab:cal2}
    %     \end{table}

\newpage
\section{Impegno orario previsto e stima dei costi totali}
    \subsection{Introduzione}
    \subsection{Stesura iniziale}

\newpage
\section{Pianificazione}
\label{sec:pianificazione}
    \subsection{Introduzione}
    \subsection{Requirements and Technology Baseline}

    % PRIMO PERIODO ----------------------------
        \subsubsection{04/11/2025 - 16/11/2025 - Primo periodo}
        
    % SECONDO PERIODO --------------------------
        \subsubsection{17/11/2025 - 30/11/2025 - Secondo periodo}

    % TERZO PERIODO --------------------------
        \subsubsection{01/12/2025 - 14/12/2025 - Terzo periodo}

    % QUARTO PERIODO --------------------------
        \subsubsection{15/12/2025 - 28/12/2025 - Quarto periodo}

    % QUINTO PERIODO --------------------------
        \subsubsection{29/12/2025 - 11/01/2026 - Quinto periodo}
    
    %\subsection{Product Baseline}
\end{document}
